\section{Interruptible Pre-Task Exploration}
\label{sec:problem}

\paragraph{Environment and memory.}
An episode takes place in an unknown environment $\mathcal{E}$.
At low-level time $t$, the agent has pose belief $x_t$, receives observation $o_t$, and maintains memory $M_t$.
Before any request is revealed, it chooses a temporally extended macro-action $a_t \in \mathcal{A}(M_t,x_t)$ with random duration $L_a$.
Examples include traversing to a frontier, approaching an uncertain object, scanning a surface, or validating a previously weak observation.
A fixed controller executes each macro-action as low-level controls $(u_{t},\ldots,u_{t+L_a-1})$.
The exploration policy never observes the realized future task.

\paragraph{Hidden task and random arrival.}
A downstream task $z\in\mathcal{Z}$ is drawn from a declared service distribution $p(z)$.
Its arrival time $\tau$ follows a possibly nonstationary hazard $h_t(s)$, the conditional probability of arrival at offset $s$ given survival until $s$.
We write $S_t(s)$ for the conditional survival function and $f_t(s)$ for the associated arrival density.
If the task arrives after low-level step $s$ of $a_t$, execution is interrupted and the pair $(M_{t+s},x_{t+s})$ is frozen.
The agent then serves $z$ with a post-arrival budget $B$.

We consider two complementary protocols.
In \textbf{recall-only}, $B=0$: the task must be answered or localized from committed evidence, isolating memory value.
In \textbf{rapid-service}, $B>0$: the task executor may navigate or acquire a small number of additional observations, so the interruption pose and reachability information also matter.
Let $U_z(M,x;B)\in[0,1]$ denote normalized downstream utility.
The objective is
\begin{equation}
\label{eq:objective}
J(\pi)=
\mathbb{E}_{\substack{\mathcal{E},\,z\sim p,\,\tau\\\pi}}
\left[
U_z(M_\tau,x_\tau;B)
-\lambda C_{0:\tau}
\right],
\end{equation}
where $C_{0:\tau}$ is pre-arrival cost.
Utility is normalized within each task family before a preregistered macro average; this prevents arbitrary scale differences between navigation and QA from determining the policy.

\paragraph{Atomic interruption semantics.}
The ordering within each low-level step is fixed:
\begin{equation}
\text{control} \rightarrow \text{observation}
\rightarrow \text{atomic memory commit}
\rightarrow \text{arrival check}.
\end{equation}
Thus the image captured on the arrival step is available, while future frames and the unexecuted suffix of a macro-action are not.
This convention is shared by every method and recorded in episode logs.

\paragraph{Information constraints.}
Before $\tau$, an admissible policy may use $M_t$, $x_t$, the service catalog and $p(z)$, and the assumed arrival model.
It may not use the sampled task, target category, arrival seed, ground-truth map, navigation mesh, shortest-path distance, or future observation.
Task-known policies are reported only as upper comparators.
All task-hidden methods share candidate actions, perception, memory capacity, low-level control, and downstream task executors.

\subsection{Why Endpoint Objectives Are Insufficient}
\label{sec:endpoint}

Let $\bar R(M,a,s)=\mathbb{E}_{z\sim p}[U_z(M\oplus\Delta M^a_s,x^a_s;B)]$
be expected service readiness after executing prefix $s$ of $a$.
An endpoint scorer compares only $\bar R(M,a,L_a)$.

\paragraph{Proposition 1 (hazard-dependent action reversal).}
There exist actions $a$ and $b$ with equal duration and equal endpoint readiness for which
(a) an endpoint-only scorer is indifferent, but
(b) the optimal action under \cref{eq:objective} changes with the arrival distribution.

\paragraph{Construction.}
Let $L_a=L_b=L$ and both endpoints satisfy one unit-valued requirement.
Action $a$ acquires the required evidence at $\epsilon<L/2$; action $b$ acquires it at $L-\epsilon$.
Their endpoint readiness is one.
For any arrival distribution concentrated on $(\epsilon,L-\epsilon)$, $a$ has strictly larger interruption utility.
Swapping the two hitting times, or placing task mass on evidence uniquely revealed early by $b$, reverses the preference while preserving terminal scores.
The appendix gives a general statement in terms of prefix-readiness integrals.
This elementary reversal is the reason \taskname{} must model evidence timing rather than attach a random termination probability to an endpoint reward.

\begin{figure}[t]
  \centering
  \placeholderfigure{3.4cm}{Controlled action-ranking reversal: two actions with matched endpoint evidence but different first-availability times under early-heavy and late-heavy hazards.}
  \caption{\textbf{Endpoint equivalence does not imply anytime equivalence.}
  The final evidence is matched, but arrival-weighted utility depends on the complete readiness curve.}
  \label{fig:reversal}
\end{figure}

